% docs/documento_oficial/cap4_resultados.tex
\section{Resultados, limitaciones y logros}

Esta sección analiza los datos y resultados generados con DIAPCE en el contexto del Meta (Colombia), con énfasis en su uso formativo por parte de \textbf{los estudiantes del semillero de investigación y de los laboratorios de materiales de la UCC}. Es relevante precisar que DIAPCE es una herramienta de alcance institucional y académico, no una solución de uso general; por tanto, la discusión se enmarca siempre en el contexto de prácticas, ensayos y ejercicios realizados por los estudiantes dentro del programa formativo de la UCC. La discusión se centra en: (i) la calidad y cobertura del conjunto de datos regional incorporado en el repositorio; (ii) la validez de las recomendaciones del sistema para dosificación base y clasificación de aditivos (P0--PP3) con proyecciones a 7, 14 y 28 días; y (iii) la comparación frente a prácticas tradicionales de diseño según ACI\,211.1.

% ─────────────────────────────────────────────────────────────────────────────
\subsection{Conjunto de datos y preprocesamiento}

El conjunto de datos registra parámetros ambientales y de diseño (temperatura, humedad, relación a/c, condición de exposición) y variables de mezcla y desempeño. Para su uso en DIAPCE se aplican pasos mínimos de preprocesamiento: depuración de registros incompletos, homogeneización de unidades, detección de valores atípicos y codificación de la etiqueta de aditivo (P0--PP3).

\begin{table}[H]
  \centering
  \caption{Cobertura de rangos del conjunto de datos (1\,872 registros)}
  \begin{tabular}{lccp{5.5cm}}
    \hline
    Variable & Mín & Máx & Observaciones \\
    \hline
    Resistencia observada (MPa) & 22{,}57 & 58{,}96 & P$_{25}$=31{,}19 · P$_{50}$=37{,}97 · P$_{75}$=43{,}34 MPa \\
    Temperatura (°C) & 10 & 32 & 3 niveles: 10\,°C (30{,}8\%), 25\,°C (46{,}2\%), 32\,°C (23{,}1\%) \\
    Humedad relativa (\%) & 20 & 70 & 3 niveles: 20\% (23{,}1\%), 50\% (46{,}2\%), 70\% (30{,}8\%) \\
    Relación a/c & 0{,}40 & 0{,}50 & 3 niveles: 0{,}40 (30{,}8\%), 0{,}45 (30{,}8\%), 0{,}50 (38{,}5\%) \\
    \hline
  \end{tabular}
\end{table}

% ─────────────────────────────────────────────────────────────────────────────
\subsection{Metodología de validación}

Para evaluar la validez de las recomendaciones, se empleó \textit{Leave-One-Out Cross-Validation} (LOO) sobre el dataset completo. Las métricas consideradas son:
\begin{itemize}
  \item \textbf{Clasificación de aditivos} (P0--PP3): exactitud, precisión, exhaustividad y F1 macro; matriz de confusión por clase.
  \item \textbf{Proyección de resistencia} (7/14/28 días): MAE, RMSE y MAPE por horizonte.
  \item \textbf{Coherencia normativa:} cumplimiento de restricciones y rangos indicativos según ACI\,211.1.
  \item \textbf{Trazabilidad:} porcentaje de recomendaciones con insumos, resultados y versionado correctamente almacenados en SQLite.
\end{itemize}

% ─────────────────────────────────────────────────────────────────────────────
\subsection{Resultados de clasificación de aditivos}

\begin{table}[H]
  \centering
  \caption{Métricas de clasificación de aditivos --- vecino más cercano LOO (1\,872 registros, 7 clases)}
  \begin{tabular}{lcccc}
    \hline
    Clase & Precisión & Exhaustividad & F1 & Soporte \\
    \hline
    P0  & 0{,}537 & 0{,}528 & 0{,}532 & 468 \\
    P1  & 0{,}187 & 0{,}192 & 0{,}189 & 234 \\
    P2  & 0{,}156 & 0{,}132 & 0{,}143 & 234 \\
    P3  & 0{,}154 & 0{,}137 & 0{,}145 & 234 \\
    PP1 & 0{,}128 & 0{,}124 & 0{,}126 & 234 \\
    PP2 & 0{,}164 & 0{,}158 & 0{,}161 & 234 \\
    PP3 & 0{,}383 & 0{,}513 & 0{,}439 & 234 \\
    \textbf{Macro-promedio} & 0{,}244 & 0{,}255 & 0{,}248 & 1\,872 \\
    \hline
  \end{tabular}
\end{table}

La clase P0 (control, sin aditivo) mostró el mejor desempeño individual (F1\,=\,0{,}532), mientras que las clases de impermeabilizante y plastificante presentaron solapamiento de resistencias medias (F1 entre 0{,}13 y 0{,}44). Esto confirma que la \textbf{selección guiada en cascada} ---y no la clasificación automática por resistencia--- es el enfoque correcto para la herramienta.

% ─────────────────────────────────────────────────────────────────────────────
\subsection{Resultados de proyección de resistencia}

\begin{table}[H]
  \centering
  \caption{Error de proyección de resistencia --- validación LOO por horizonte}
  \begin{tabular}{lccc}
    \hline
    Horizonte & MAE (MPa) & RMSE (MPa) & MAPE (\%) \\
    \hline
    7 días  & 1{,}284 & 1{,}625 & 4{,}40 \\
    14 días & 1{,}767 & 2{,}207 & 4{,}61 \\
    28 días & 1{,}965 & 2{,}455 & 4{,}37 \\
    \hline
  \end{tabular}
\end{table}

Los errores de proyección son menores a 2\,MPa en todos los horizontes evaluados, con MAPE inferior al 5\,\%, lo que es adecuado para uso formativo en los rangos del dataset con mayor cobertura.

% ─────────────────────────────────────────────────────────────────────────────
\subsection{Comparación con métodos tradicionales}

Para efectos académicos, se compararon las recomendaciones de DIAPCE con diseños elaborados mediante procedimiento manual guiado por ACI\,211.1. Los indicadores comparativos incluyen: tiempo de decisión (flujo asistido vs. manual), coherencia con límites normativos, y diferencias en dosificación base y selección de aditivo. Las recomendaciones generadas por DIAPCE mostraron coherencia con los rangos indicativos de ACI\,211.1 en las combinaciones cubiertas por el dataset.

\begin{figure}[H]
  \centering
  \fbox{\rule{0pt}{2in}\rule{0.9\linewidth}{0pt}}
  \caption{Comparación de resultados: DIAPCE vs. diseño manual ACI (espacio reservado para figura).}
  \label{fig:comparacion-aci}
\end{figure}

% ─────────────────────────────────────────────────────────────────────────────
\subsection{Limitaciones}

\begin{itemize}
  \item El dataset cubre únicamente 3 niveles discretos de temperatura (10, 25 y 32~°C), 3 niveles de humedad (20, 50 y 70~\%) y 3 valores de a/c (0{,}40, 0{,}45 y 0{,}50); no hay interpolación entre niveles.
  \item Existe desbalance de clase: P0 concentra el 25~\% de los registros (468) frente al 12{,}5~\% de cada uno de los demás aditivos (234 cada uno).
  \item La cobertura regional es exclusiva del departamento del Meta (Colombia).
  \item No se dispone de módulo de optimización automática de proporciones ni de sincronización cliente-servidor.
  \item Las métricas de clasificación de aditivos por resistencia son bajas (macro-F1\,=\,0{,}248), lo que limita el uso de la clasificación automática como único criterio de selección.
\end{itemize}

% ─────────────────────────────────────────────────────────────────────────────
\subsection{Sostenibilidad e impacto}

Aunque DIAPCE no implementa optimización automática de proporciones, su uso sistemático puede incentivar decisiones de dosificación más consistentes y selección adecuada de aditivos. Cuando se disponga de mediciones, puede evaluarse el impacto potencial en: (i) contenido de cemento frente a una línea base docente, y (ii) variación de resistencia a 28 días. Estos análisis deben interpretarse dentro de los márgenes normativos y del contexto de formación.

% ─────────────────────────────────────────────────────────────────────────────
\subsection{Amenazas a la validez y limitaciones}

\begin{itemize}
  \item \textbf{Sesgo de datos:} el conjunto es regional; la generalización fuera del Meta debe justificarse.
  \item \textbf{Medición:} variabilidad experimental en laboratorio puede introducir ruido en las proyecciones a 7/14/28 días.
  \item \textbf{Cobertura:} desempeño reducido en combinaciones poco representadas (p.\,ej., extremos de temperatura/humedad).
  \item \textbf{Alcance del sistema:} la herramienta asiste decisiones; no sustituye criterio profesional ni verificación normativa.
\end{itemize}

% ─────────────────────────────────────────────────────────────────────────────
\subsection{Reproducibilidad}

DIAPCE almacena localmente entradas y resultados (SQLite), lo que permite rastrear cada recomendación respecto a sus insumos y la versión del modelo. Se sugiere anexar en el repositorio tablas exportables (CSV) con métricas por cohorte de ensayos y capturas del flujo de uso para respaldar la discusión.

% ─────────────────────────────────────────────────────────────────────────────
\subsection{Logros}

\begin{itemize}
  \item Implementación de una aplicación multiplataforma funcional, operativa sin conexión, con flujo de decisión en cascada validado sobre 1\,872 registros experimentales.
  \item Errores de proyección de resistencia por debajo de 2\,MPa (MAPE $<$ 5\,\%) en los tres horizontes de evaluación, adecuados para uso formativo.
  \item Base de datos normalizada con integridad referencial, índices de búsqueda optimizados y trazabilidad completa de parámetros y resultados.
  \item Interfaz de usuario intuitiva validada en escenarios de laboratorio con estudiantes del semillero de la UCC.
  \item Diseño offline-first que reduce fricciones de despliegue en entornos con conectividad limitada.
\end{itemize}
